\chapter{参数高效微调}
\label{lab:42_peft}

\noindent\textbf{配套文件：}\path{notebook/42_peft.ipynb}\par
\noindent\textbf{教材关联：}\bookxrefshort{ch:parameter-efficient-finetuning}、\bookxrefshort{ch:efficient-finetuning}。

\section*{实验目标}
验证低秩增量的参数范围、合并等价性和行为验收。

\section*{准备及定位}
先阅读本册实验总则，确认Notebook中的依赖与资产单元。原理计算、库接口迁移和真实模型训练可能具有不同资源要求，应分别记录设备、版本及运行范围。

源码定位可从下列函数开始；应结合前置数据与配置单元阅读。
\begin{itemize}
\item \texttt{my\_parameter\_budget}
\item \texttt{my\_lora\_story\_color\_limit}
\item \texttt{my\_draw\_lora\_story\_cells}
\end{itemize}

\section*{操作及观察}
\begin{enumerate}
\item 计算低秩参数预算，核对只有预期模块可训练。
\item 训练低秩增量并比较合并前、合并后、拆分及重载的固定输入输出。
\item 进入Qwen3身份任务时冻结聊天模板、回答掩码、训练记录与评估探针。
\item 记录适配器制品与底座身份，对照训练前后行为及重载后的固定探针。
\end{enumerate}

\section*{结果判读及提交}
提交参数清单、数值误差和分切片行为结果。身份回答改变不等于全部任务能力改善；LoRA和量化底座训练应分别记录精度与执行条件。

提交时附上所执行单元范围、环境与制品身份、原始结果及失败说明。将未运行项目明确列出，不能用Notebook中已有输出替代本次执行证据。

相关方法的原始定义与适用前提见\cite{huLoRALowRankAdaptation2022}；本实验的运行结果仍须按实际配置单独验证。


框架实践把参数注入、监督目标和适配器制品连接起来，使用Transformers负责模型与分词器、PEFT负责LoRA、TRL负责SFT训练\cite{HuggingfaceTransformers2018,huggingfacePEFT2023,huggingface2026trlframework}。配套入口为\path{book/workbook/examples/train_adapter.py}，从仓库根目录运行，与前面的Notebook共享机制目标，不依赖其内存状态。

准备一个已下载并固定修订的因果聊天模型目录，以及来源明确的训练与验证JSONL文件。入口要求支持BF16的CUDA设备；所需显存取决于基座、长度和激活，不能仅由LoRA参数量估计。可先使用能够在单卡容纳的小型指令模型验证完整流程，再扩大规模。模型目录应包含权重、配置、分词器、聊天模板及相应特殊词元信息。

接口核查基线为TRL 1.13.0与PEFT 0.20.0；训练依赖与推理引擎分别建环境。先按目标设备安装兼容的PyTorch，再安装并检查训练依赖，记录解析得到的完整版本集合。以下操作会安装软件，应在正式运行实验时执行。
\begin{verbatim}
python -m venv .venv-train
source .venv-train/bin/activate
python -m pip install "trl==1.13.0" "peft==0.20.0" \
  "accelerate==1.15.0" datasets
python -m pip check
python -m pip freeze > /absolute/run/train-packages.txt
\end{verbatim}
\path{/absolute/run}须替换为已建立的实验记录目录。依赖解析成功只说明声明的包约束得到满足，设备算子、模型加载和反向计算仍需运行验证。重现实验时复用经验证的锁定环境，升级则重新核对数据处理与训练行为。

每行包含\texttt{prompt}与\texttt{completion}两列，值为角色消息数组。下面仅展示记录格式，训练集应使用经过审校且有使用权限的数据；验证集按文档、来源或任务划分，避免同源近重复样本泄漏。
\begin{verbatim}
{"prompt":[{"role":"user","content":"请解释梯度累积。"}],"completion":[{"role":"assistant","content":"梯度累积在多个微批次上累计梯度，再执行一次参数更新。"}]}
\end{verbatim}
入口要求最后一条提示来自用户，且待学习部分恰好是一条助手回答。它检查训练与验证提示的完全重复、模板后的长度及首批监督目标；近重复和语义泄漏仍须在数据准备阶段检查。采用\texttt{completion\_only\_loss=True}，监督回答部分；\texttt{packing=False}让首次验收保留每条记录的边界。

正式更新前，应从\texttt{trainer.get\_train\_dataloader()}取得一个批次，对照\texttt{input\_ids}与\texttt{labels}解码非\texttt{-100}位置，确认回答及结束标记受到监督，提示没有被错误计入。源码中的非空检查只能发现没有训练目标的情况，不能替代这项语义检查。EOS参数应取聊天模板规定的回答结束词元，不应仅凭词元名称猜测。

把路径替换为实际的绝对路径，并从模型模板核实\texttt{EOS\_TOKEN}。下面的100步用于完整流程验收，不代表已达到某项质量目标；步数、学习率和验证频率应根据数据与验证曲线调整。
\begin{verbatim}
python book/workbook/examples/train_adapter.py \
  --model /absolute/model-snapshot \
  --train /absolute/data/train.jsonl \
  --eval /absolute/data/validation.jsonl \
  --output /absolute/run/lora-01 \
  --eos-token "$EOS_TOKEN" \
  --steps 100 --max-length 1024 --accumulation 8
\end{verbatim}
固定单卡微批次为1，累积8步；LoRA秩为16，缩放参数为32，目标为PEFT识别的线性层。先核对打印或保存的可训练参数名称，确认输出头和额外保存层是否符合任务要求。入口拒绝向已有非空输出目录开始新训练；恢复时必须显式指定检查点。

\path{final/}保存适配器和分词器，\path{checkpoint-*}用于训练恢复，\path{experiment.json}记录参数及数据哈希。另行保存基座的不可变修订或文件哈希、驱动与库版本、显存和时间记录。仅有本地路径不能唯一标识基座内容。

在新的进程中加载同一基座，再通过\texttt{PeftModel.from\_pretrained}附加\path{final/}。使用相同模板、精度和\texttt{eval()}模式，比较保存前后固定输入的logits或损失，并记录容差；随后在保留测试集上比较微调前后任务质量与能力退化。需要合并时，在可容纳的浮点基座上执行\texttt{merge\_and\_unload}并保存新的完整目录，独立验证后再交给推理引擎。量化基座、合并和重新量化的误差应分别测量。

LLaMA-Factory与Axolotl适合把重复的数据处理、训练和导出组织成配置工作流\cite{llamafactory2026framework,axolotl2026framework}。迁移时以同一数据与制品为输入，逐项映射下表。字段拼写按所固定的版本查阅，不能只复制不同版本的配置文件。
\begin{center}
\begin{tabularx}{\linewidth}{p{30mm}XX}
\toprule 比较对象 & 必须保持的语义 & 验收证据\\\midrule
数据与模板 & 样本划分、角色、EOS和回答监督范围 & 同一样本的词元与标签。\\
适配参数 & 目标模块、秩、缩放及额外保存层 & 可训练参数名称与数量。\\
优化过程 & 有效全局批量、更新次数、精度与学习率 & 配置展开结果与训练日志。\\
制品 & 基座身份、适配器与分词器绑定 & 独立重载和固定探针结果。\\
\bottomrule
\end{tabularx}
\end{center}
TRL中的DPO或GRPO训练器对应另一类数据与反馈契约。偏好训练需要同一提示下的候选偏好，策略训练还要规定采样和奖励；把训练器名字替换掉，不能使SFT数据自动成为有效的偏好或强化学习数据。

提交依赖记录、数据划分、监督位置样例、可训练参数、验证损失和独立重载结果。损失不变时先检查标签、学习率与参数是否获得梯度；输出异常结束时核对模板和EOS；显存不足时分别分析权重、激活、长度与批量。不得用单条身份回答变化替代保留集评估。
